\chapter{Architettura}

\section{Panoramica}

L'architettura di Open Data Chunker segue un pattern \textbf{layered} con separazione netta delle responsabilità. Il sistema è composto da quattro layer principali, illustrati nel diagramma seguente.

\begin{figure}[h]
\centering
\begin{tikzpicture}[node distance=1.5cm, scale=0.9, transform shape]
    
    % Stile database migliorato
    \tikzstyle{dbbox} = [rectangle, rounded corners=3pt, minimum width=3.5cm, minimum height=1.2cm, text centered, draw=black, fill=accentgreen!25, font=\small]
    
    % Input Layer
    \node[component, fill=warnorange!30, minimum width=4cm] (xml) {Input Layer\\{\small XML Files}};
    
    % Processing Layer
    \node[component, fill=primaryblue!20, minimum width=4cm, right=2cm of xml] (process) {Processing Layer\\{\small lxml + Sanitizer}};
    
    % Storage Layer - Box stile database
    \node[dbbox, right=2cm of process] (storage) {Storage Layer\\{\small Apache Parquet}};
    \draw[black, thick] ([yshift=-3pt]storage.north west) -- ([yshift=-3pt]storage.north east);
    \draw[black, thick] ([yshift=-8pt]storage.north west) -- ([yshift=-8pt]storage.north east);
    
    % Query Layer
    \node[component, fill=primaryblue!40, minimum width=4cm, below=2cm of storage] (query) {Query Layer\\{\small DuckDB + Polars}};
    
    % CLI
    \node[component, fill=gray!20, minimum width=3cm, below=2cm of process] (cli) {CLI\\{\small Click}};
    
    % Arrows
    \draw[arrow] (xml) -- node[above, font=\small] {stream} (process);
    \draw[arrow] (process) -- node[above, font=\small] {batch write} (storage);
    \draw[arrow] (storage) -- (query);
    \draw[arrow] (cli) -- (process);
    \draw[arrow] (cli) -- (query);
    
\end{tikzpicture}
\caption{Architettura a layer di Open Data Chunker}
\label{fig:architecture}
\end{figure}

\section{Layer di Input}

Il layer di input gestisce i file XML grezzi provenienti dall'RNA Open Data. Caratteristiche:

\begin{itemize}
    \item Supporto per file singoli o directory ricorsive
    \item Discovery automatico tramite glob pattern \texttt{*.xml}
    \item Nessun caricamento in memoria dell'intero file (streaming)
\end{itemize}

\section{Layer di Processing}

Il core della pipeline, implementato in \texttt{src/parser.py}. Composto da:

\begin{table}[h]
\centering
\begin{tabular}{@{}lp{8cm}@{}}
\toprule
\textbf{Componente} & \textbf{Responsabilità} \\
\midrule
\texttt{CleanFileInputStream} & Wrapper che filtra caratteri XML non validi in tempo reale \\
\texttt{iterparse()} & Parser SAX-like per elaborazione streaming \\
\texttt{ProcessPoolExecutor} & Orchestratore del parallelismo multi-processo \\
\texttt{flush\_batches()} & Writer batch su Parquet con partizionamento \\
\bottomrule
\end{tabular}
\caption{Componenti del Processing Layer}
\end{table}

\section{Layer di Storage}

Output in formato \textbf{Apache Parquet} con le seguenti caratteristiche:

\begin{itemize}
    \item \textbf{Columnar storage}: compressione ottimale e predicate pushdown
    \item \textbf{Hive partitioning}: directory \texttt{ANNO=YYYY} per partition pruning
    \item \textbf{Schema tipizzato}: definizioni PyArrow in \texttt{src/models.py}
\end{itemize}

\begin{tipbox}[Vantaggio: Partition Pruning]
Le query filtrate per anno leggono solo le partizioni rilevanti, riducendo I/O fino al 90\% su dataset multi-anno.
\end{tipbox}

\section{Layer di Query}

Dual-engine per flessibilità:

\begin{description}
    \item[DuckDB] Motore SQL analitico in-process. Ideale per aggregazioni e join complessi via CLI.
    \item[Polars] DataFrame library lazy-first. Utilizzato per export streaming memory-efficient.
\end{description}

\section{Flusso dei Dati}

\begin{enumerate}
    \item L'utente invoca \texttt{parse --input data/} via CLI
    \item Il CLI discovery tutti i file \texttt{*.xml} nella directory
    \item I file vengono distribuiti tra N worker processes
    \item Ogni worker:
    \begin{enumerate}
        \item Apre il file tramite \texttt{CleanFileInputStream}
        \item Itera gli elementi \texttt{<AIUTO>} con \texttt{iterparse()}
        \item Estrae i campi in batch da 10.000 record
        \item Scrive il batch su Parquet partizionato
    \end{enumerate}
    \item Al termine, statistiche aggregate vengono loggite
\end{enumerate}
